\documentclass[fleqn]{article}
\usepackage{amsmath, amssymb, amsthm}
\usepackage{fullpage}
\usepackage{listings}
\usepackage{stix}
\usepackage[utf8]{inputenc}
\lstset{ basicstyle=\ttfamily, }

\begin{document}

\section*{Excercise 3.3 1}

\subsection*{(a): Write a program to compute $f(x) = \begin{cases}0 & x = 0 \\ 1 & x \neq 0\end{cases}$}

Let $g(x, y)$ be the function in part (c).
Then $f(x) = g(x, 0)$, and the program in (c) computes $x=0$.

\subsection*{(b): Write a program to compute $f(x) = 5$}

\begin{lstlisting}
Z(1)
S(1)
S(1)
S(1)
S(1)
S(1)
\end{lstlisting}

\subsection*{(c) Write a program to compute $f(x, y) = \begin{cases}0 & x = y \\ 1 & x \neq y\end{cases}$}

\begin{lstlisting}
J(1, 2, 3)
S(3)
T(3, 1)
\end{lstlisting}

\section*{Excercise 3.3 3}

\subsection*{A program without J}

Assume that no T instructions are in the program.
Let $M$ be the longest sequence of S(1) instructions such that no Z(1) occurs after the final S(1) instruction.
Let $m$ be the length of that sequence.
If no Z(1) instuctions occur before the sequence $M$, then Register 1 will contain the value $x+m$.
Otherwise, Register 1 will contain the value $m$.

Relaxing the assumption that no T instructions occur in the program, observe that T insructions can be of the form $T(1, j)$ (form a), $T(j, 1)$ (form b) or $T(j, k)$ (form c).
By the logic above, let $m_a$ be the number of non-interrupted $S(1)$ instructions.
Then after a $T(1, j)$ (form a) instruction, register $j$ will hold either $x+m_a$ or $m_a$.

For a $T(j, 1)$ (form b) instruction, by the above logic after a $T(1, j)$ (form a) instruction, the register j will hold $x+m_a$ or $m_a$.
By the above logic with a change of index to $S(j)$, an additional $m_b$ may be accumulated in register $k$, bringing its value to $x+m_a+m_b$.
Let $m = m_a + m_b$.
Then after a form b instruction with no following $Z(1)$ instructions, register 1 will hold the value $m$ or $x+m$.

By induction, this process can be extended to n $T(j, k)$ (form c) instructions, so long as the chain is ended by a form b instruction, and possibly started with a form a instruction.

\section*{Excercise 3.3 4}

\subsection*{Show that T(m, n) is not strictly necessary.}

Let I be the index of the T(m, n) instruction. Modify the program as such:

First, for every J instruction after the T instruction, increase the destination value by 4.

Second, replace the T(m, n) instruction with the following four instructions.

\begin{lstlisting}
Z(n)
J(m, n, I+4)
S(n)
J(m, n, I+2)
\end{lstlisting}

\section*{Excercise 4.3}

\subsection*{(a) Show that $x<y$ is computable.}

Let $f(x, y) = \begin{cases}0&x=y\\1&x<y\\2&x>y\end{cases}$. See email dated Feb 6.

\subsection*{(b) Show that $x\neq3$ is computable.}

Let $f(x, y)$ be the function in 4.3 a.

Let $g(x) = \begin{cases}0&x = 3\\>0&x\neq3\end{cases} = f(x, 3)$.

See 4.3 a.

\section*{Excercise 5.2 1}
\subsection*{Show $2x$ on $\mathbb{Z}$ is decidable}

Example 5.1 shows $\alpha: \mathbb{X}\to\mathbb{N}$ and $\alpha^-1: \mathbb{N}\to\mathbb{Z}$.
Let $f^*(x) = \alpha(2x)\alpha^-1$. $f(n) = 2x$ is show in example ..., so $f^*(x)$ is decidable.
Expanded, $f^*(x) =
\begin{cases}2x\\2x+1\end{cases} f(x)\alpha^-1 =
\begin{cases}4x\\4x+2\end{cases} \alpha^-1 =
\begin{cases}2x & \text{   } x \text{ even}\\-2(x-1) & \text {   } x \text{ odd}\end{cases}
$

\end{document}